\section*{Абстракт}
\addcontentsline{toc}{section}{Абстракт}

В тази курсова работа е разработена работеща платформа за персонализирано обучение, използваща MongoDB като нерелационна база данни. Проектът демонстрира правилно структуриране на данни чрез три колекции (потребители, курсове, напредък), изпълнение на CRUD операции с различни видове филтриране, агрегиращи заявки за анализ на данни и управление на достъпа чрез ролево-базиран контрол (RBAC). Реализирана е REST API с Express.js и интерактивен уеб интерфейс, който позволява лесно тестване на всички функционалности. Базата данни съдържа над 30 реалистични записа, включително вложени документи и референции между колекциите. Проектът показва предимствата на нерелационните бази данни при работа с гъвкави и йерархични структури от данни в сферата на електронното обучение.
